\chapter{Containerd}
\label{ch:containerd}

Containerd is the container runtime at the heart of Veilbox. It manages
the full container lifecycle — pulling images, creating snapshots,
starting processes, and monitoring execution. This chapter explains its
architecture, configuration, and integration with the rest of the system.

\section{Containerd Architecture}

Containerd is organized into several layers:

\begin{figure}[H]
\centering
\begin{tikzpicture}[node distance=0.8cm]
\node[highlight] (nerdctl) {nerdctl (CLI)};
\node[box, below of=nerdctl] (grpc) {gRPC API};
\node[box, below of=grpc] (svc) {Services Layer};
\node[box, below left of=svc, xshift=-1.5cm] (images) {Images};
\node[box, below of=svc] (snap) {Snapshots};
\node[box, below right of=svc, xshift=1.5cm] (cont) {Containers};
\node[box, below of=snap] (runc) {runc (OCI)};
\draw[arrow] (nerdctl) -- (grpc);
\draw[arrow] (grpc) -- (svc);
\draw[arrow] (svc) -- (images);
\draw[arrow] (svc) -- (snap);
\draw[arrow] (svc) -- (cont);
\draw[arrow] (snap) -- (runc);
\draw[arrow] (cont) -- (runc);
\end{tikzpicture}
\caption{Containerd layered architecture. nerdctl communicates via gRPC
with the containerd daemon, which manages images, snapshots, and
container processes through runc.}
\label{fig:containerd}
\end{figure}

\subsection{gRPC API Layer}

Containerd exposes a gRPC API on a Unix domain socket at
\texttt{/run/containerd/containerd.sock}. The API is organized into
several services:

\begin{itemize}[nosep]
  \item \textbf{Namespace service}: Create, list, remove namespaces.
  - \textbf{Image service}: Pull, push, list, tag, remove images.
  - \textbf{Container service}: Create, start, stop, delete containers.
  - \textbf{Task service}: Attach to container processes, read logs.
  - \textbf{Snapshotter service}: Manage filesystem snapshots for image
        layers.
  - \textbf{Diff service}: Create and apply layer diffs.
\end{itemize}

\subsection{Client Libraries}

Clients can interact with containerd via:
\begin{itemize}[nosep]
  \item \textbf{ctr}: The containerd-native CLI (minimal, debugging).
  - \textbf{nerdctl}: Docker-compatible CLI (user-friendly).
  - \textbf{Go client}: \texttt{github.com/containerd/containerd} Go
        package for programmatic access.
  - \textbf{containerd-shim}: Shim process that sits between containerd
        and runc, enabling live-upgrade of containerd without stopping
        containers.
\end{itemize}

\section{Image Management}

\subsection{Image Pull Flow}

When a user runs \texttt{nerdctl pull alpine:latest}:

\begin{enumerate}[nosep]
  \item nerdctl sends a \texttt{PullImage} gRPC request to containerd.
  \item containerd resolves the image reference (e.g.,
        \texttt{docker.io/library/alpine:latest}) through configured
        registries.
  \item containerd authenticates with the registry (if credentials are
        configured).
  \item containerd downloads the image manifest (JSON) which lists the
        configuration and layer blobs.
  \item For each layer blob, containerd verifies the digest (SHA256),
        downloads the compressed blob, and decompresses it.
  \item The snapshotter creates a new snapshot for each layer.
  \item The image is registered in containerd's content store and
        metadata database.
\end{enumerate}

\subsection{Image Storage}

Images are stored in two locations:

\begin{itemize}[nosep]
  \item \textbf{Content store}: \texttt{/mnt/state/containerd/blobs/}.
        Blobs are stored by content hash (SHA256). Each blob is a
        subdirectory named by the first two hex characters of the hash.
  \item \textbf{Metadata}: \texttt{/mnt/state/containerd/metadata/}.
        BoltDB database tracking image references, layer metadata, and
        container state.
\end{itemize}

\section{Container Lifecycle}

\subsection{Creating a Container}

\begin{lstlisting}
nerdctl run -d --name web nginx:alpine
\end{lstlisting}

Steps:
\begin{enumerate}[nosep]
  \item containerd creates a new container object in the metadata store.
  \item The snapshotter creates a writable snapshot on top of the image
        layers (using OverlayFS).
  \item containerd generates an OCI specification (config.json) that
        describes the container's rootfs, mounts, environment variables,
        command, resource limits, and security settings.
  \item containerd starts the containerd-shim process, passing the
        container ID and bundle path.
  \item The shim launches runc with the OCI bundle, creating the
        container's namespaces, cgroups, and process.
\end{enumerate}

\subsection{Running Containers}

\begin{lstlisting}
nerdctl ps                    # List running containers
nerdctl exec -it web sh       # Execute command in container
nerdctl logs web              # View container logs
nerdctl stats                 # Resource usage
nerdctl stop web              # Graceful stop (SIGTERM)
nerdctl rm web                # Remove container
\end{lstlisting}

\paragraph{Docker CLI compatibility.}
The \texttt{docker} command is available as a symlink to \texttt{nerdctl}.
Every \texttt{nerdctl} example in this chapter also works with \texttt{docker}:
\texttt{docker run}, \texttt{docker ps}, \texttt{docker exec}, etc.
No Docker daemon runs inside the VM — \texttt{nerdctl} (and therefore
\texttt{docker}) communicates directly with containerd via gRPC.

\section{Containerd Configuration}

\begin{lstlisting}
# /etc/containerd/config.toml
version = 2
root = "/mnt/state/containerd"
state = "/run/containerd"

[grpc]
  address = "/run/containerd/containerd.sock"

[plugins]
  [plugins."io.containerd.grpc.v1.cri"]
    sandbox_image = "registry.k8s.io/pause:3.9"
  [plugins."io.containerd.snapshotter.v1.overlayfs"]
    root_path = "/mnt/state/containerd/snapshots"
\end{lstlisting}

\section{Integration with Veilbox}

Containerd is started in rcS:

\begin{lstlisting}[language=bash]
containerd -c /etc/containerd/config.toml &
sleep 2  # Wait for gRPC socket
nerdctl pull alpine:latest &
\end{lstlisting}

The gRPC socket at \texttt{/run/containerd/containerd.sock} is available
for local clients. No network port is exposed (containerd listens on a
Unix socket only).

\section{Containerd Namespaces}

Containerd supports multiple namespaces for tenant isolation:

\begin{lstlisting}
# List namespaces
ctr namespaces ls

# Create a namespace
ctr namespaces create tenant1

# Run container in namespace
nerdctl --namespace tenant1 run alpine echo hello

# View containers by namespace
nerdctl --namespace tenant1 ps -a
nerdctl --namespace k8s.io ps -a  # Kubernetes namespace
\end{lstlisting}

\section{Image Management Internals}

\subsection{Image Pull}

When pulling an image, containerd performs these steps:

\begin{lstlisting}
# Trace image pull
ctr images pull docker.io/library/alpine:latest

# Show image metadata
ctr images ls -q

# Export image to tar
ctr images export alpine.tar docker.io/library/alpine:latest

# Import image from tar
ctr images import alpine.tar
\end{lstlisting}

\subsection{Content Addressable Storage}

Images are stored using content-addressable storage (CAS). Each blob
is identified by its SHA256 digest:

\begin{lstlisting}
# Content store location
ls /mnt/state/containerd/blobs/sha256/

# View blob metadata
cat /mnt/state/containerd/blobs/sha256/<hash>

# Check content store size
du -sh /mnt/state/containerd/
\end{lstlisting}

\section{Snapshotter Plugins}

Containerd supports multiple snapshotter backends:

\begin{longtable}{|p{3.5cm}|p{4cm}|p{4cm}|}
\hline
\textbf{Snapshotter} & \textbf{Description} & \textbf{Performance} \\
\hline
overlayfs & Native OverlayFS (recommended) & Best \\
native & Copy-on-write via hard links & Worst \\
devmapper & Device-mapper thin provisioning & Medium \\
btrfs & Btrfs snapshots & Good \\
fuse-overlayfs & FUSE-based OverlayFS & Slow \\
nydus & Nydus snapshotter (lazy pull) & Fast pull \\
\hline
\end{longtable}

Veilbox uses overlayfs by default, which provides the best performance
for container workloads.

\section{Containerd Event Monitoring}

Containerd emits events for monitoring and automation:

\begin{lstlisting}
# Stream all events (in real time)
ctr events

# Filter events by topic
ctr events | grep "image/pull"

# List recent events (from log)
grep containerd /var/log/messages | tail -20
\end{lstlisting}

\section{Resource Requirements}

Containerd in Veilbox requires:

\begin{itemize}[nosep]
  \item \textbf{Disk space}: Minimum 256 MB free on state partition for
        image pulls. Recommended: 1+ GB for multiple images.
  \item \textbf{Memory}: Approximately 30 MB for the containerd daemon
        itself. Container memory usage depends on the workload.
  \item \textbf{File descriptors}: Containerd may use 100+ file
        descriptors when managing multiple containers.
\end{itemize}

Container resource limits are set per-container:

\begin{lstlisting}
# Memory limit
nerdctl run --memory=256m nginx:alpine

# CPU limit
nerdctl run --cpus=0.5 nginx:alpine

# PIDs limit
nerdctl run --pids-limit=100 nginx:alpine

# Block I/O limit
nerdctl run --blkio-weight=500 nginx:alpine
\end{lstlisting}

\section{Logging Driver}

Containerd supports configurable logging:

\begin{lstlisting}
# Configuring log rotation in config.toml
[plugins."io.containerd.grpc.v1.cri".containerd]
  [plugins."io.containerd.grpc.v1.cri".containerd.default_runtime]
    runtime_type = "io.containerd.runc.v2"

[plugins."io.containerd.grpc.v1.cri".containerd.runtimes.runc]
  runtime_type = "io.containerd.runc.v2"

[plugins."io.containerd.grpc.v1".cni]
  bin_dir = "/opt/cni/bin"
  conf_dir = "/etc/cni/net.d"
\end{lstlisting}
